\section{验证矩阵与 Go/No-Go 判定}

在正式替换生产训练路径前，必须完成以下判定实验。任何关键项失败都应先回到理论或数据设计阶段，而不是通过继续堆网络层数掩盖问题。

\subsection{Gate 1：电流二次型恒等式}
随机采样 $(a,g,u)$，比较原始 reduced EM/Joule 计算与
\begin{equation}
q_j^{\mathrm{quad}}=\zeta^\top G_j(a,g)\zeta.
\end{equation}
要求误差达到由浮点和 EM solver tolerance 决定的数值等价水平。若该测试失败，必须首先检查激励参数化、复数共轭、offset 或 loss operator 定义；禁止进入神经网络阶段。

\subsection{Gate 2：$G$ 输出可压缩性}
在代表性 $(a,g)$ 数据集上对 $y=\operatorname{vec}_{\mathrm{sym}}G$ 做 SVD。记录：
\begin{itemize}
  \item singular-value decay；
  \item $K=16,32,64,128,\ldots$ 的 held-out RMS 与 maximum reconstruction error；
  \item 对 $q$ 的最坏电流方向误差。
\end{itemize}
若低 rank 压缩不成立，不否定总体架构；退化策略为增加 $K$ 或完整输出 $y$。只有当完整 $G$ 映射本身也表现出不可学习的非平滑/高频结构时才重新审视模型。

\subsection{Gate 3：热状态有效维数}
分析
\begin{equation}
\frac{\partial y}{\partial a}
\quad\text{或}\quad
\frac{\partial q}{\partial a}
\end{equation}
的统计 singular spectrum/active subspace。该结果用于判断未来是否值得对 198 维热状态做输入压缩，但首版不以此作为架构成立条件。

\subsection{Gate 4：稳定性区域}
在完整目标参数域的高覆盖样本上计算真实
\begin{equation}
J_F=-M^{-1}K+M^{-1}\frac{\partial q}{\partial a}
\end{equation}
及合适能量范数下的 logarithmic norm。输出至少包括：
\begin{itemize}
  \item 最坏观测 one-sided Lipschitz 常数；
  \item 是否存在可信的 uniform negative margin；
  \item 哪些 geometry/current/temperature 区域接近失稳。
\end{itemize}
只有在有充分证据时才启用 contractive hard constraint。

\subsection{Gate 5：神经映射学习能力}
使用普通 residual MLP 训练并在冻结 test set 上分别报告：
\begin{align}
E_G &= \norm{G_\theta-G},\\
E_q &= \norm{q_\theta-q},\\
E_F &= \norm{F_\theta-F}.
\end{align}
必须同时报告 RMS、95/99 percentile 和 maximum，避免平均误差掩盖 hard regions。

\subsection{Gate 6：轨迹与输出验证}
选择覆盖中心、边界、高温、高电流和几何极端的初值/工况，使用以下参考层级：
\begin{enumerate}
  \item neural ODE/ETD vs. Radau + same EM-ROM；
  \item selected audit cases：Radau + full sparse EM；
  \item 稳态：neural steady solve 后用 real reduced/full physics residual audit。
\end{enumerate}
报告：
\begin{itemize}
  \item maximum coordinate error；
  \item maximum nodal-temperature error；
  \item maximum temperature 的误差；
  \item impedance/region loss error（若查询）；
  \item 长时间 drift；
  \item time-step refinement convergence。
\end{itemize}

\subsection{Gate 7：性能验证}
必须在同一硬件和相同物理问题上记录：
\begin{itemize}
  \item 物理标签生成吞吐量；
  \item NN 训练 wall time 与 GPU 利用率；
  \item 单次 vector-field evaluation；
  \item 1s/10s/100s trajectory query；
  \item 大批量参数扫描吞吐；
  \item 内存与模型文件尺寸。
\end{itemize}
不以理论 FLOP 推算代替真实 benchmark。

\subsection{生产替换条件}
只有当以下条件同时成立时，新架构才替换当前 production surrogate：
\begin{enumerate}
  \item Gate 1 精确通过；
  \item frozen test set 上 $q/F$ 误差达到预设工程预算；
  \item trajectory/output 误差满足目标；
  \item long-time 行为无系统漂移；
  \item 训练流程在可接受资源预算内稳定复现；
  \item 推理性能满足生产要求；
  \item 模型保存/加载/版本签名和数据版本可复现。
\end{enumerate}

若新架构尚未满足替换条件，旧系统只作为临时 baseline 使用；不再继续扩展其网络复杂度作为主要解决方案。
